\chapter{Năng lượng chân không, lạm phát và các thăng giáng nguyên thủy}

\section{Năng lượng chân không và không gian de Sitter}

Trong không gian phẳng, phương trình Friedmann là
\begin{equation}
\left(\frac{\dot A}{A}\right)^2=\frac{8\pi G}{3}\rho.
\end{equation}
Đối với vật chất, \(\rho\) giảm theo \(A^{-3}\). Đối với bức xạ, \(\rho\) giảm theo \(A^{-4}\). Đối với năng lượng chân không, nó không bị pha loãng:
\begin{equation}
\rho=\rho_0=\text{constant}.
\end{equation}
Nếu vũ trụ không chứa gì ngoài năng lượng chân không này, vế phải là một hằng số. Viết hằng số đó là \(H^2\),
\begin{equation}
H^2=\frac{8\pi G}{3}\rho_0,
\qquad
H=\sqrt{\frac{8\pi G}{3}\rho_0}.
\end{equation}
Khi đó hằng số Hubble thực sự là không đổi:
\begin{equation}
\frac{\dot A}{A}=H.
\end{equation}
Do đó
\begin{equation}
\dot A=HA,
\end{equation}
và nghiệm có dạng hàm mũ:
\begin{equation}
A(t)=C e^{Ht}.
\end{equation}
Khoảng cách giữa các thiên hà xa không liên kết hấp dẫn tăng theo hàm mũ. Đây là không gian de Sitter.

Trong các tọa độ này, hình học là
\begin{equation}
d\tau^2=dt^2-e^{2Ht}\left(dx^2+dy^2+dz^2\right).
\end{equation}
Năng lượng chân không ở đây là dương. Trong không gian phẳng, nó buộc phải như vậy, vì vế trái của phương trình Friedmann là không âm. Năng lượng chân không âm chỉ có thể xuất hiện khi có kể đến các hạng độ cong, và đó không phải là trường hợp đang xét ở đây.

Einstein đã đưa vào hằng số vũ trụ học vì một mục đích khác: cân bằng vật chất thông thường và giữ cho vũ trụ tĩnh.

\section{Vũ trụ hiện nay đang gia tốc}

Không gian de Sitter cũng là xu hướng ở thời muộn của vũ trụ hiện nay. Trước đó, thừa số tỉ lệ tuân theo dạng ứng xử do vật chất chi phối
\begin{equation}
A(t)\propto t^{2/3}.
\end{equation}
Vài tỉ năm trước, đường cong đã uốn lên trên và bắt đầu tăng nhanh hơn. Nó vẫn chưa phải là một hàm mũ hoàn hảo, vì các dạng năng lượng khác vẫn còn hiện diện, nhưng xu hướng đang tiến về dạng đó.

Gia tốc ở đây có nghĩa là
\begin{equation}
\ddot A>0.
\end{equation}
Đây là một đặc điểm đã được quan sát của sự giãn nở ở thời muộn.

Bây giờ hãy xét chân trời. Định luật Hubble cho biết
\begin{equation}
v=HD.
\end{equation}
Tại khoảng cách mà tốc độ lùi xa bằng tốc độ ánh sáng,
\begin{equation}
D=\frac{c}{H}.
\end{equation}
Tương đương, có thể viết lại điều này thành \(HD=c\).\editorialnote{Đây là cách diễn đạt lại mang tính biên tập về mặt đại số của hệ thức được phát biểu bằng lời \(D=c/H\). Dòng trên cùng trên bảng bên phải trong Hình~\ref{fig:lecture07-horizon} chỉ đọc được một phần.}

\begin{figure}[ht]
\centering
\includegraphics[width=0.62\textwidth]{lecture_07_figure_02.png}
\caption{Quan hệ chân trời vũ trụ \(D=c/H\) trong giãn nở de Sitter.}
\label{fig:lecture07-horizon}
\lectureframe{00:30:37}
\end{figure}

Khoảng cách này có cỡ từ mười đến mười lăm tỉ năm ánh sáng, gần bằng về mặt số học với tuổi của vũ trụ khi đại lượng sau được quy đổi thành một khoảng cách bằng \(c\).

\begin{classroomqa}
\lecturetimestamp{00:14:10}
\audiencequestion{Chẳng phải sự gần bằng nhau giữa thang chân trời và tuổi của vũ trụ chỉ là điều xảy ra nếu đường tiếp tuyến tình cờ đi qua gốc tọa độ sao?}
\lecturerresponse{Điều đó sẽ đúng chính xác đối với \(A(t)\propto t\). Nói chung thì không đúng. Trong không gian de Sitter, \(H\) giữ nguyên còn tuổi thì tiếp tục tăng, nên sự gần bằng nhau hiện nay không phải là một đồng nhất thức toán học.}
\end{classroomqa}

\section{Ý nghĩa của năng lượng chân không}

Ở đây năng lượng chân không được mô tả bằng ngôn ngữ của cơ học lượng tử thông thường. Đối với một dao động tử điều hòa, năng lượng có dạng sơ lược là
\begin{equation}
E\sim p^2+x^2.
\end{equation}
Theo cơ học cổ điển, trạng thái cơ bản là điểm cân bằng,
\begin{equation}
x=0,\qquad p=0,\qquad E=0.
\end{equation}
Cơ học lượng tử không cho phép đồng thời cả \(x\) và \(p\) đều được cố định chính xác bằng không. Nguyên lý bất định buộc phải có các thăng giáng, nên năng lượng trạng thái cơ bản là dương. Đối với dao động tử, ta được
\begin{equation}
E_0=\frac{1}{2}\hbar\omega.
\end{equation}

Ý này sau đó được chuyển sang các trường. Vũ trụ đầy những dao động tử điều hòa: các dao động của trường điện từ, và tương tự như vậy đối với mọi trường khác. Mỗi trường đều có năng lượng điểm không. Đối với trường điện từ, mật độ năng lượng có dạng sơ lược là
\begin{equation}
u_{\mathrm{EM}}\sim E^2+B^2.
\end{equation}
Ngay cả khi không có photon thực nào hiện diện, vẫn còn các dao động điểm không không thể triệt tiêu. Đó là một cách giải thích năng lượng chân không. Tính chất duy nhất cần dùng ở đây là năng lượng này có mặt và không bị pha loãng khi vũ trụ giãn nở.

\begin{classroomqa}
\lecturetimestamp{00:16:15}
\audiencequestion{Năng lượng chân không này là gì, và vì sao vẫn còn bất cứ thứ gì nếu vật chất và bức xạ bị loại bỏ?}
\lecturerresponse{Vì chân không là trạng thái cơ bản của các trường lượng tử, và một trạng thái cơ bản lượng tử không nhất thiết phải có năng lượng bằng không. Các dao động tử điều hòa có năng lượng trạng thái cơ bản dương, còn các trường thì ứng xử như những tập hợp dao động tử. Vì vậy, ngay cả khi không có hạt thực nào hiện diện, vẫn còn năng lượng điểm không.}
\end{classroomqa}

Đây không phải là bức xạ nền ba độ. Bức xạ nền vi sóng là bức xạ thông thường. Mật độ năng lượng của nó giảm vì các photon bị pha loãng và mỗi photon đều bị dịch đỏ. Năng lượng chân không là phần năng lượng điểm không còn lại sau khi các photon thực đã bị loại bỏ.

\section{Ý nghĩa của \(H^{-1}\)}

Bây giờ quay lại mối liên hệ giữa nghịch đảo của hằng số Hubble và tuổi của vũ trụ. Nếu
\begin{equation}
A(t)\propto t,
\end{equation}
thì
\begin{equation}
H=\frac{\dot A}{A}=\frac{1}{t}.
\end{equation}
Tổng quát hơn, nếu
\begin{equation}
A(t)\propto t^p,
\end{equation}
thì
\begin{equation}
H=\frac{p}{t}.
\end{equation}
Đối với trường hợp vật chất chi phối \(p=2/3\),
\begin{equation}
H=\frac{2}{3t}.
\end{equation}
Vì vậy, trong các mô hình vũ trụ học dạng lũy thừa, \(H^{-1}\) có cùng bậc độ lớn với tuổi, nhưng không chính là tuổi.

Không gian de Sitter thì khác. Ở đó
\begin{equation}
A(t)=C e^{\gamma t},
\end{equation}
và do đó
\begin{equation}
H=\frac{\dot A}{A}=\gamma,
\end{equation}
một hằng số thực sự. Ở đây \(H^{-1}\) không phải là tuổi của vũ trụ. Nó là hằng số thời gian của sự tăng trưởng mũ: khoảng thời gian trong đó hệ số tỉ lệ được nhân với \(e\). Trong phần trình bày bằng lời, người nói nhiều lần dùng ``gấp đôi'' rồi tự đính chính; phát biểu chính xác là được nhân với \(e\). Ngày nay hằng số thời gian đó có cỡ mười tỉ năm.

\begin{classroomqa}
\lecturetimestamp{00:23:46}
\audiencequestion{Thầy/cô có thể nói lại điểm về tuổi của vũ trụ và nghịch đảo của hằng số Hubble được không?}
\lecturerresponse{Với các luật lũy thừa, ta có \(H=p/t\), nên \(H^{-1}\) xấp xỉ tuổi của vũ trụ, sai khác chỉ bởi một hệ số cỡ một. Nhưng trong không gian de Sitter, \(H\) là hằng số trong khi tuổi vẫn tiếp tục tăng, nên \(H^{-1}\) chỉ là hằng số thời gian của sự giãn nở theo hàm mũ, không phải tuổi.}
\end{classroomqa}

\section{Năng lượng và chân trời vũ trụ}

Nếu mật độ năng lượng chân không giữ không đổi trong khi vũ trụ giãn nở, thì tự nhiên ta sẽ hỏi liệu tổng năng lượng chân không có đang tăng lên hay không. Câu trả lời ở đây chỉ giới hạn trong một phạm vi nhất định. Trong thuyết tương đối rộng, việc hạch toán năng lượng bao gồm cả một thành phần hấp dẫn, nên khái niệm tổng năng lượng tinh tế hơn so với trong cơ học sơ cấp. Điểm có ý nghĩa thực tế là phần gắn với chân trời.

Trong không gian de Sitter chính xác, chân trời vũ trụ có kích thước cố định,
\begin{equation}
D=\frac{c}{H}
=\frac{c}{\sqrt{8\pi G\rho_0/3}}.
\end{equation}
Do đó thể tích bên trong chân trời ấy là cố định. Năng lượng chân không trong vùng đó là
\begin{equation}
E_{\mathrm{hor}}=\frac{4\pi}{3}D^3\rho_0.
\end{equation}
\editorialnote{Đoạn này cô đọng phép tính bằng lời tại 00:31:06--00:32:23: lấy kích thước chân trời, tính thể tích hình cầu, rồi nhân với \(\rho_0\).}
Vì cả \(D\) lẫn \(\rho_0\) đều là hằng số, đại lượng này không thay đổi theo thời gian. Đối với phần vũ trụ mà một quan sát viên cho trước có thể quan sát được, năng lượng chân không mà họ có thể tiếp cận vẫn cố định trong không gian de Sitter chính xác.

Các vật thể cũng có thể đi qua chân trời vũ trụ gần như theo đúng nghĩa mà chúng đi qua chân trời lỗ đen. Một khi đã ở bên ngoài, chúng không còn nằm trong miền liên hệ nhân quả với chúng ta.

\begin{classroomqa}
\lecturetimestamp{00:29:27}
\audiencequestion{Nếu mật độ năng lượng chân không giữ không đổi trong khi vũ trụ giãn nở, thì phần năng lượng tăng thêm này đến từ đâu?}
\lecturerresponse{Bài giảng trả lời bằng ngôn ngữ của thuyết tương đối rộng, chứ không phải ngôn ngữ cái hộp kiểu Newton. Trong không gian de Sitter chính xác, kích thước chân trời giữ cố định, nên năng lượng chân không bên trong vùng mà một quan sát viên có thể quan sát được về nguyên tắc cũng giữ cố định.}
\end{classroomqa}

\section{Vì sao lạm phát được đưa ra}

Sau đó, chính cơ chế de Sitter ấy được đẩy lùi về quá khứ xa xôi, nhưng với một hằng số thời gian ngắn hơn rất nhiều, cỡ
\begin{equation}
H^{-1}\sim 10^{-32}\,\mathrm{s}.
\end{equation}
Đó là ý tưởng của lạm phát: trước bất cứ điều gì có thể quan sát trực tiếp, vũ trụ đã trải qua một giai đoạn giãn nở theo hàm mũ.

Động lực cơ bản là độ phẳng và độ trơn nhẵn của vũ trụ. Các lát cắt không gian có vẻ rất phẳng. Ở các thang lớn, vũ trụ cũng rất đồng nhất. Đáng chú ý hơn nữa, mặt tán xạ cuối cùng đã trơn nhẵn sẵn rồi, nên sự làm trơn đó không thể chỉ diễn ra trong giai đoạn lịch sử khả kiến về sau.

Một quả bóng xẹp thì nhăn nheo. Bơm nó đủ căng lên thì bất kỳ vùng nhìn thấy nào cũng trở nên nhẵn. Phép loại suy với nho khô hay mận khô cũng nói lên cùng một điểm: một vật nhăn nheo, khi bị kéo giãn với một hệ số cực lớn, sẽ trông nhẵn trên mọi vùng nhìn thấy được. Lạm phát được cho là đã làm đúng như vậy.

Điều này cũng giải đáp bài toán chân trời theo cách dùng ở đây. Những vùng cách xa nhau trên mặt tán xạ cuối cùng đều quan sát được đối với chúng ta, nhưng ánh sáng đã không có đủ thời gian để truyền từ vùng này sang vùng kia trong phần lịch sử khả kiến về sau. Vì vậy, mật độ và nhiệt độ gần như đồng đều của chúng gợi ý một tiền sử trong đó một vùng nhỏ hơn đã bị kéo giãn một cách hết sức mạnh mẽ.

Lạm phát có một ưu điểm hai mặt. Nó xóa bỏ độ nhạy đối với các điều kiện ban đầu, và cũng vì cùng lý do đó, nó xóa bỏ phần lớn thông tin về việc vũ trụ đã khởi đầu như thế nào.

Vấn đề độ phẳng lại xuất hiện ở đây dưới dạng sắc nét hơn. Về mặt quan sát, tham số độ cong \(k\) có vẻ bằng không, ở mức khoảng một phần trăm hoặc tốt hơn. Hình học trong vùng khả kiến có vẻ Euclid. Lạm phát giải thích điều này bằng cách làm cho bán kính độ cong lớn hơn rất nhiều so với vùng chúng ta có thể thấy.

\begin{classroomqa}
\lecturetimestamp{00:40:05}
\audiencequestion{Có phải ta đang giả sử có cân bằng trước lạm phát không?}
\lecturerresponse{Không. Ý chính của lạm phát gần như là điều ngược lại. Bất kỳ cấu trúc nào có ở đó từ trước đều bị kéo giãn mạnh đến mức vùng khả kiến không còn lưu giữ thông tin về các chi tiết trước đó.}
\end{classroomqa}

\section{Một trường vô hướng và thế năng của nó}

Để giải thích vì sao năng lượng chân không hiệu dụng ở vũ trụ sơ khai có thể đã lớn hơn hiện nay rất nhiều, ta đưa vào một trường vô hướng. Không giống điện trường và từ trường, vốn là các vectơ, trường vô hướng không có hướng. Nó đơn giản chỉ có một giá trị tại mỗi điểm:
\begin{equation}
\phi=\phi(x,y,z,t).
\end{equation}
Trường này là giả thuyết. Nó không được đồng nhất với bất kỳ trường nào đã biết trong tự nhiên, và các lượng tử của nó cũng không được đồng nhất với bất kỳ hạt nào đã biết.

Phần năng lượng của nó quan trọng ở đây là thế năng,
\begin{equation}
V(\phi).
\end{equation}
Cụm ``thế năng'' được hiểu là năng lượng gắn với chính giá trị của trường, chứ không gắn với đạo hàm theo thời gian của trường. Các giá trị khác nhau của \(\phi\) có thể có những năng lượng khác nhau.

\begin{figure}[ht]
\centering
\includegraphics[width=0.68\textwidth]{lecture_07_figure_03.png}
\caption{Thế vô hướng định tính \(V(\phi)\) cho thấy rằng các giá trị trường khác nhau mang các thế năng khác nhau.}
\label{fig:lecture07-generic-potential}
\lectureframe{00:47:08}
\end{figure}

\begin{classroomqa}
\lecturetimestamp{00:46:34}
\audiencequestion{Trục ngang trong hình \(V(\phi)\) là gì?}
\lecturerresponse{Trục ngang là giá trị trường \(\phi\), còn trục đứng là thế năng \(V\).}
\end{classroomqa}

Nếu trường không đổi theo không gian và thời gian, hoặc gần như không đổi trên miền đang xét, thì \(V(\phi)\) ứng xử đúng hệt như năng lượng chân không. Nó không bị pha loãng. Vì thế các công thức de Sitter trước đó được dùng lại với \(\rho_0\) được thay bằng \(V(\phi)\). Cụ thể,
\begin{equation}
H=\sqrt{\frac{8\pi G}{3}V(\phi)}.
\end{equation}

Nếu \(\phi\) thay đổi theo vị trí, các vùng khác nhau sẽ lạm phát với các hằng số thời gian khác nhau. Nhưng nếu sự kéo giãn đủ lớn, bất kỳ miền khả kiến nào cũng trở nên gần như đồng nhất.

\begin{classroomqa}
\lecturetimestamp{01:05:11}
\audiencequestion{Ở đây \(\phi\) có đang đóng vai trò như một tham số thời gian không?}
\lecturerresponse{Trong thảo luận này thì có, theo nghĩa hạn chế là trong khi trường lăn xuống thế năng thì nó là một hàm đơn điệu theo thời gian.}
\end{classroomqa}

\section{Một cao nguyên, một vách đá và sự tái nung nóng}

Thế năng cần cho câu chuyện này trở nên rất phẳng, dù không hoàn toàn phẳng, rồi đến mép một vách đá, giảm mạnh, không hẳn xuống tới không, và lại tăng lên.

\begin{figure}[ht]
\centering
\includegraphics[width=0.68\textwidth]{lecture_07_figure_04.png}
\caption{Thế lạm phát với một cao nguyên cao, một đoạn giảm rất dốc, và một cực tiểu thấp hơn.}
\label{fig:lecture07-inflationary-potential}
\lectureframe{01:02:30}
\end{figure}

Dạng này được sắp đặt rất gượng ép, nhưng nó được đưa vào vì câu chuyện vũ trụ học dường như đòi hỏi một cái gì đó giống như thế.

Trường bắt đầu ở vị trí cao trên cao nguyên. Vì độ dốc nhỏ và độ nhớt hiệu dụng gắn với sự giãn nở lại lớn, \(\phi\) chỉ thay đổi rất chậm. Một hòn đá rơi chậm qua mật ong gần với hình ảnh được muốn gợi ra hơn là một vật trượt trên mặt bàn. Vũ trụ giãn nở càng nhanh thì độ nhớt này càng lớn.

\begin{classroomqa}
\lecturetimestamp{00:59:07}
\audiencequestion{Điều gì khiến \(\phi\) dịch chuyển?}
\lecturerresponse{Độ dốc của thế đẩy nó. Độ dốc nhỏ, còn độ nhớt vũ trụ học lớn, nên sự dịch chuyển là chậm chứ không nhanh.}
\end{classroomqa}

Trong khi \(\phi\) nhích từng chút một dọc theo cao nguyên, \(V(\phi)\) thay đổi rất ít. Vì thế vũ trụ lạm phát gần như thể năng lượng chân không là hằng số. Hằng số thời gian trong giai đoạn này vào cỡ
\begin{equation}
H^{-1}\sim 10^{-32}\,\mathrm{s}.
\end{equation}
Một ước lượng thô bằng lời cho độ cao của cao nguyên vào khoảng \(10^{-14}\) theo đơn vị Planck.\editorialnote{Đây là ước lượng ``phỏng đoán tốt nhất'' được nêu bằng lời tại 01:14:38--01:14:47, chứ không phải một giá trị được suy ra trong bài giảng.}

Lạm phát kéo dài bao lâu? Điều đó chưa biết, vì chưa biết bề rộng và độ nông của cao nguyên. Điều đã biết là vũ trụ đã được nhân thêm một hệ số \(e\) ít nhất khoảng 50 lần, và 60 là con số thường được nêu.\editorialnote{Đây là bản biên tập diễn đạt lại phần thảo luận bằng lời về e-folding tại 01:14:05--01:19:27.} Đây là các e-folding: mỗi e-fold là một lần nhân với \(e\).

Khi trường tới mép, thế năng biến đổi nhanh chóng. Trường rơi xuống, thế năng của nó trước hết chuyển thành động năng rồi thành nhiệt và các hạt, và vũ trụ trở nên nóng. Tên gọi ``tái nung nóng'' hơi gây hiểu lầm trong mô tả này, vì mọi năng lượng nhiệt sớm hơn đã bị pha loãng mất trong thời kỳ lạm phát.

Cực tiểu thấp hơn không ở mức không. Giá trị nhỏ còn lại ở đó chính là giá trị gắn với sự giãn nở mũ chậm hiện nay.

\begin{classroomqa}
\lecturetimestamp{01:19:48}
\audiencequestion{Liệu cao nguyên có thể phẳng hoàn toàn, còn các nhiễu loạn bằng cách nào đó hất trường khỏi mép hay không?}
\lecturerresponse{Không. Nó không thể hoàn toàn phẳng. Nó phải có một độ nghiêng nào đó. Nếu quá phẳng, vũ trụ sau tái nung nóng sẽ trở nên quá không đồng nhất.}
\end{classroomqa}

``Vụ Nổ Lớn'' được dùng ở đây theo hai nghĩa gần nhau. Nó có thể chỉ chính sự kiện tái nung nóng. Nó cũng có thể chỉ giai đoạn sớm nhất có thể thấy trực tiếp, tức mặt tán xạ cuối cùng. Điểm chắc chắn là mọi thứ thực sự được quan sát đều nằm rất xa bên dưới cao nguyên lạm phát.

\begin{classroomqa}
\lecturetimestamp{01:45:42}
\audiencequestion{Bản thân sự tái nung nóng có phải là Vụ Nổ Lớn không?}
\lecturerresponse{Nếu thích thì có. Nhưng ở thời điểm đó vũ trụ vẫn cực kỳ nóng và mờ đục. Mặt tán xạ cuối cùng đến muộn hơn rất nhiều, sau khi vũ trụ đã nguội đủ để trở nên trong suốt.}
\end{classroomqa}

\section{Thăng giáng lượng tử và mầm mống của các thiên hà}

Lạm phát làm cho vũ trụ khả kiến trở nên rất đồng nhất, nhưng không hoàn toàn đồng nhất. Nếu không có các thăng giáng lượng tử, trường sẽ trở nên hoàn toàn trơn tru. Nhưng với cơ học lượng tử, nó vẫn biến thiên nhẹ từ nơi này sang nơi khác.

Hình minh họa trung tâm biểu diễn không gian thông thường theo phương thẳng đứng và giá trị trường \(\phi\) theo phương ngang. Nếu trường hoàn toàn đồng nhất, thì ở một thời điểm cho trước mọi điểm trong không gian sẽ có cùng \(\phi\), và mọi điểm sẽ cùng trôi về phía mép. Khi đó mọi vùng sẽ rơi qua vách tại cùng một thời điểm. Các thăng giáng lượng tử phá vỡ sự đồng bộ chính xác đó.

Có hai hình ảnh làm rõ ý này. Một là một gợn nhỏ trong một trường vốn đồng nhất ở những chỗ khác. Hình kia là một hàng lính khoác tay nhau đi về phía một vách đá. Nếu họ thẳng hàng hoàn toàn, tất cả sẽ rơi cùng nhau. Nếu một số hơi đi trước và một số hơi tụt lại sau, họ sẽ đi qua mép vào những thời điểm khác nhau.

Bây giờ đến hiệu ứng quan trọng. Một vùng rơi qua mép sớm hơn thì tái nung nóng sớm hơn. Đến khi một vùng muộn hơn rơi qua, vùng sớm hơn đã có thời gian giãn nở và nguội đi phần nào. Sự khác biệt đó tạo ra các biến thiên của mật độ năng lượng sau tái nung nóng.

\begin{classroomqa}
\lecturetimestamp{01:34:00}
\audiencequestion{Các vùng được tái nung nóng muộn hơn nóng hơn hay lạnh hơn các vùng sớm hơn?}
\lecturerresponse{Những vùng sớm hơn đã có thời gian giãn nở và nguội đi. Vì thế những vùng muộn hơn nóng hơn khi chúng rơi qua, và sự khác biệt đó góp phần tạo nên sự không đồng nhất của mật độ.}
\end{classroomqa}

Hiệu ứng được quan sát là nhỏ:
\begin{equation}
\frac{\delta\rho}{\rho}\sim 10^{-5},
\end{equation}
xấp xỉ một phần trong một trăm nghìn. Đây là các thăng giáng mật độ nguyên thủy, không phải năng lượng tối. Con số này ở đây không được suy ra chỉ từ riêng phác họa định tính; nó được dùng để ràng buộc cao nguyên có thể cao đến đâu và phẳng đến mức nào.

Chúng cũng là điều cần thiết. Không có chúng sẽ không có các thiên hà. Sau tái nung nóng, hấp dẫn khuếch đại mẫu hình ban đầu. Các vùng có mật độ cao hút vật chất từ các vùng có mật độ thấp. Các vùng có mật độ cao trở nên càng có mật độ cao hơn, các vùng có mật độ thấp càng có mật độ thấp hơn, và kết quả là cấu trúc về sau của các thiên hà và sao.

Đây không chỉ là một câu chuyện định tính. Ta có thể tính các thăng giáng lượng tử của một trường trong một vũ trụ đang giãn nở, theo dõi chúng qua giai đoạn tái nung nóng, rồi so sánh kết quả với mẫu hình thăng giáng được phục dựng từ phân bố các thiên hà. Sự phù hợp là rất tốt.

Cùng thảo luận ấy cũng dẫn đến một điểm về entropy. Sự kết tụ do hấp dẫn trông có vẻ phản trực giác, nhưng nó không vi phạm định luật hai. Trong cơ học thống kê, hấp dẫn khá khác thường: nó có xu hướng làm cho mọi thứ kết tụ lại, và sự kết tụ đó làm tăng entropy chứ không làm giảm nó.

\begin{classroomqa}
\lecturetimestamp{01:40:24}
\audiencequestion{Bức xạ nền vi sóng được thể hiện trong toàn bộ điều này như thế nào, khi \(\phi\) không phải là một trường điện từ?}
\lecturerresponse{Mẫu mật độ nguyên thủy được in dấu lên bề mặt tán xạ cuối cùng. Tái nung nóng chuyển năng lượng của trường thành các hạt thông thường và photon, và về sau bức xạ nền vi sóng mang dấu ấn của những thăng giáng nguyên thủy đó.}
\end{classroomqa}

\section{Sóng âm và bước tiếp theo}

Cuối cùng, phần thảo luận chuyển ngắn gọn sang sóng âm. Chúng diễn ra trước sự kết tụ do hấp dẫn về sau. Các thăng giáng nguyên thủy được biến đổi thành các dao động âm, và những dao động ấy để lại những cụm có kích thước vật lý đã biết trên bầu trời.

Điều đó quan trọng vì một kích thước vật lý đã biết ở một khoảng cách đã biết cho ta một thước đo chuẩn. Một khi đo được góc mà nó chắn trên bề mặt tán xạ cuối cùng, ta có thể dùng hình học để xác định không gian là phẳng, cong dương hay cong âm. Các chi tiết được hoãn lại.

\begin{classroomqa}
\lecturetimestamp{01:47:19}
\audiencequestion{Còn sóng âm thì sao?}
\lecturerresponse{Chúng cung cấp một thước đo chuẩn trên bầu trời. Kích thước vật lý đã biết của chúng, cùng với kích thước góc được quan sát trên bề mặt tán xạ cuối cùng, là điều cho phép ta đo hình học của không gian. Các chi tiết được dời sang bài giảng tiếp theo.}
\end{classroomqa}